\documentclass[10pt,a4paper]{article}
\usepackage[T1]{fontenc}
\usepackage[utf8]{inputenc}
\usepackage{lmodern}
\usepackage[margin=20mm]{geometry}
\usepackage{amsmath,amssymb,array,longtable,booktabs,calc,enumitem}
\usepackage{xurl,xcolor,mdframed}
\usepackage{float}
\usepackage{graphicx}
\usepackage[unicode,hidelinks]{hyperref}
\usepackage{kotex}
\setlength{\parindent}{0pt}\setlength{\parskip}{5pt}
\setlength{\emergencystretch}{3em}\setlength{\tabcolsep}{4pt}
\renewcommand{\arraystretch}{1.12}\allowdisplaybreaks
\newcolumntype{L}[1]{>{\raggedright\arraybackslash}p{#1}}
\DeclareMathOperator{\CVaR}{CVaR}
\DeclareMathOperator{\logit}{logit}
\begin{document}

%%%%%%%%%%

\section{2단계 구조의 ER-DRO 모형 (CVaR형 ambiguity set)}

데이터: $(x^k,p^k)$, $k=1,\dots,N$. 예를 들어 로보택시가 운영된 도시--기간 단위로 관측된 context $x^k$(무차량 가구 비율)와 로보택시 선택확률 $p^k$이다. 회귀로 $\hat f$와 로짓 척도의 잔차 $\hat\varepsilon^k=\logit(p^k)-(\hat\beta_0+\hat\beta_1x^k)$를 얻고(\S\ref{sec:data}), 후보 도시 $c$의 context $x_c$에 잔차를 하나씩 붙여 시나리오 $D_c^s$를 생성한다. 잔차 하나가 시나리오 하나에 대응하므로 시나리오 집합 $S$는 잔차 인덱스 집합이고 $|S|=N$이다.
아래에서는 모형을 1단계 문제 (1), 시나리오별 2단계 문제 (3), 그리고 실제로 풀이하는 확장형(extensive form) (4)로 나누어 쓴다. 1단계 변수는 $(y,q)$이고, 2단계 변수 $(O^s,U^s)$는 시나리오 $s$마다 분리된다.
차량 1대당 배치·운영비는 별도의 항으로 두지 않는다. 모형의 모든 항이 통행 단위로 표현되므로 대당 비용 $K_c$는 용량 1통행당 비용 $K_c/\kappa_c$로 환산되어 $c_c^+$, $c_c^-$에 흡수되며(\S\ref{sec:absorb}), 이를 독립된 항으로 두어도 최적해는 바뀌지 않는다.

\subsection{1단계 문제}

\begin{subequations}\label{eq:first}
\begin{align}
\min_{y,q}\quad
 & \underbrace{\sum_{c\in C}F_c y_c}_{\text{진입 비용}}
 + \underbrace{\mathcal R_\zeta(q)}_{\text{시나리오 페널티의 CVaR}} \label{eq:1a}\\
\text{s.t.}\quad
 & 0\le q_c\le \bar q_c\,y_c, && \forall c\in C, \label{eq:1b}\\
 & \sum_{c\in C} q_c\le Q, \label{eq:1c}\\
 & y_c\in\{0,1\},\ q_c\in\mathbb Z_{\ge 0}, && \forall c\in C. \label{eq:1d}
\end{align}
\end{subequations}
여기서 $\mathcal R_\zeta(q)$는 2단계 비용 $H(q,D^s)$의 시나리오 분포에 대한 CVaR이며, Rockafellar--Uryasev의 threshold 표현으로
\begin{equation}\label{eq:cvar}
\mathcal R_\zeta(q)=\min_{\tau\in\mathbb R}\Big\{\tau+\frac{1}{|S|(1-\zeta)}\sum_{s\in S}\big(H(q,D^s)-\tau\big)^+\Big\}
\end{equation}
로 쓴다. $\zeta=0$이면 $\mathcal R_0(q)=\frac1{|S|}\sum_s H(q,D^s)$가 되어 ER-SAA(표본평균)로 환원된다. Context $x_c$는 시나리오 값 $D_c^s=M_c p_c^s$에만 들어가고 최적화 구조는 바뀌지 않는다.

\subsection{2단계 문제 (시나리오별 recourse)}

1단계 결정 $q$와 시나리오 $s$의 수요 $D^s=(D_c^s)_{c\in C}$가 주어졌을 때 2단계 비용은
\begin{subequations}\label{eq:second}
\begin{align}
H(q,D^s)=\min_{O^s,U^s}\quad
 & \sum_{c\in C}\big(c_c^+O_c^s+c_c^-U_c^s\big) \label{eq:3a}\\
\text{s.t.}\quad
 & \kappa_c q_c-D_c^s=O_c^s-U_c^s, && \forall c\in C, \label{eq:3b}\\
 & O_c^s,\ U_c^s\ge 0, && \forall c\in C. \label{eq:3c}
\end{align}
\end{subequations}
(3)은 $c$별로도 분리되며, 최적해는 닫힌 형태 $O_c^s=(\kappa_c q_c-D_c^s)^+$, $U_c^s=(D_c^s-\kappa_c q_c)^+$이다. 따라서 $H(q,D^s)=\sum_c\big[c_c^+(\kappa_c q_c-D_c^s)^+ + c_c^-(D_c^s-\kappa_c q_c)^+\big]$는 $q$에 대해 볼록인 조각별 선형 함수이고, 이것이 2단계 분해(L-shaped) 적용의 근거가 된다.
이 2단계 비용은 신문판매원(newsvendor) 문제의 과잉·부족 비용과 같은 구조이며, 1단계 결정 $q$는 (2)의 CVaR을 통해서만 2단계와 연결된다.

\subsection{확장형 (extensive form): 실제로 풀이하는 MILP}

(2)의 $(\cdot)^+$를 변수 $w^s$로 선형화하고 (3)을 (1)에 합치면 다음 단일 MILP를 얻는다. (4b)--(4c), (4g)--(4h)는 시나리오별로 분리되는 2단계 제약이고, (4d)--(4f)는 1단계 제약이다.
\begin{subequations}\label{eq:ext}
\begin{align}
\min_{y,q,O,U,\tau,w}\quad
 & \sum_{c\in C}F_c y_c+\tau+\frac{1}{|S|(1-\zeta)}\sum_{s\in S}w^s \label{eq:4a}\\
\text{s.t.}\quad
 & w^s\ge\sum_{c\in C}\big(c_c^+O_c^s+c_c^-U_c^s\big)-\tau, && \forall s\in S, \label{eq:4b}\\
 & \kappa_c q_c-D_c^s=O_c^s-U_c^s, && \forall c\in C,\ s\in S, \label{eq:4c}\\
 & 0\le q_c\le\bar q_c\,y_c, && \forall c\in C, \label{eq:4d}\\
 & \sum_{c\in C}q_c\le Q, \label{eq:4e}\\
 & y_c\in\{0,1\},\ q_c\in\mathbb Z_{\ge0}, && \forall c\in C, \label{eq:4f}\\
 & O_c^s,\ U_c^s\ge 0, && \forall c\in C,\ s\in S, \label{eq:4g}\\
 & w^s\ge 0,\ \ \tau\in\mathbb R, && \forall s\in S. \label{eq:4h}
\end{align}
\end{subequations}

\subsection{변형: 진입 비용을 예산 제약으로}

진입 가능한 도시 수(또는 진입 예산)가 실제 제약이라면, (4a)의 $\sum_c F_c y_c$를 목적함수에서 제거하고 다음 제약을 추가한다.
\begin{equation}\label{eq:budget}
\sum_{c\in C}F_c y_c\le B \qquad\text{또는}\qquad \sum_{c\in C}y_c\le B.
\end{equation}
이 변형에서는 목적함수가 $\tau+\frac{1}{|S|(1-\zeta)}\sum_s w^s$, 즉 2단계 페널티의 CVaR만 남는다. 이 변형은 $y$의 조합을 직접 제한하므로 2단계 분해 시 초기 master 문제가 자명해($y=q=0$) 반복이 길어지는 문제를 완화한다.

\subsection{대당 비용의 흡수 : $K_c$ 항 빼는 이유}\label{sec:absorb}

차량 1대당 비용 $K_c$(배치·운영비)를 목적함수에 $\sum_c K_c q_c$로 두더라도, (4c)에 의해 모든 $s$에서 $\kappa_c q_c=D_c^s+O_c^s-U_c^s$이므로 $\zeta=0$일 때
\begin{equation}\label{eq:absorb}
\sum_{c}K_c q_c+\frac1{|S|}\sum_{s}\sum_{c}\big(c_c^+O_c^s+c_c^-U_c^s\big)
=\underbrace{\frac1{|S|}\sum_s\sum_c\frac{K_c}{\kappa_c}D_c^s}_{\text{결정과 무관한 상수}}
+\frac1{|S|}\sum_{s}\sum_{c}\Big[\Big(c_c^++\frac{K_c}{\kappa_c}\Big)O_c^s+\Big(c_c^--\frac{K_c}{\kappa_c}\Big)U_c^s\Big]
\end{equation}
이 성립한다. 즉 $K_c$는 $c_c^+\leftarrow c_c^++K_c/\kappa_c$, $c_c^-\leftarrow c_c^--K_c/\kappa_c$의 재정의와 동치이며($\zeta>0$에서는 $(K_c/\kappa_c)D_c^s$가 CVaR 안에 들어가 근사), 독립된 결정 레버가 아니다. 또한 $K_c$가 도시 간에 같고 (4e)가 등호로 묶이면 $\sum_c K_c q_c=KQ$로 상수이다. 따라서 본 연구는 $K_c$를 별도 항으로 두지 않고, 대당 비용을 용량 1통행당으로 환산한 값을 $c_c^+$, $c_c^-$의 정의에 포함한다 생각하고 비율만 따진다(ex 1:5).

{\footnotesize
\begin{longtable}{@{}L{0.34\textwidth}L{0.12\textwidth}L{0.50\textwidth}@{}}
\toprule 제약 & 단계 & 의미\\\midrule\endhead\bottomrule\endlastfoot
$0\le q_c\le\bar q_c y_c$ \eqref{eq:4d} & 1단계 & 도시에 진입하지 않으면 차량 0대, 진입하면 도시별 상한까지 배치 가능\\
$\sum_c q_c\le Q$ \eqref{eq:4e} & 1단계 & 총 차량 대수 제한\\
$\sum_c F_c y_c\le B$ \eqref{eq:budget} & 1단계 (변형) & 진입 예산(또는 진입 도시 수) 제한. 목적함수의 $\sum_c F_c y_c$를 대체\\
$\kappa_c q_c-D_c^s=O_c^s-U_c^s$ \eqref{eq:4c} & 2단계 ($s$별 분리) & 서비스 용량과 로보택시 수요의 차이 $=$ 잉여 용량 $-$ 미충족 수요\\
$w^s\ge\sum_c(c_c^+O_c^s+c_c^-U_c^s)-\tau$, $w^s\ge0$ \eqref{eq:4b},\eqref{eq:4h} & 2단계 ($s$별 분리) & 시나리오 $s$의 페널티가 threshold $\tau$를 초과하는 양. 최소화 과정에서 $w^s=(\text{penalty}_s-\tau)^+$가 된다\\
\end{longtable}}

\section{표기}
\subsection{Contextual 및 DRO 추가 표기}
{\footnotesize
\begin{longtable}{@{}L{0.14\textwidth}L{0.45\textwidth}L{0.35\textwidth}@{}}
\toprule 기호 & 의미 & 종류 및 단위\\\midrule\endhead\bottomrule\endlastfoot
$(x^k,p^k)$ & 관측 $k$의 context(무차량 가구 비율)와 로보택시 선택확률 & 데이터; $k=1,\dots,N$\\
$x_c$ & 후보 도시 $c$의 특성(공변량) 벡터. 현 단계에서는 무차량 가구 비율 하나 & 입력; Kannan의 새 관측 $x$에 대응\\
$\hat f$ & 추정된 회귀함수; $\hat f(x)\approx\mathbb E[p\mid x]$. 로짓 $\hat f(x)=\sigma(\hat\beta_0+\hat\beta_1x)$, $\sigma(z)=1/(1+e^{-z})$ & 추정 함수\\
$\hat\varepsilon^k$ & 잔차 $\logit(p^k)-(\hat\beta_0+\hat\beta_1x^k)$ & 유도량; 로짓 척도\\
$\zeta$ & CVaR형 ambiguity set의 반경(꼬리 비율). $1-\zeta$는 최악 그룹으로 남기는 시나리오 비율 & 입력; $[0,1)$. $\zeta=0$이면 ER-SAA\\
$\tau$ & 페널티 threshold. 결정변수이며 최적해에서 ``$|S|(1-\zeta)$번째로 큰 페널티'' & 결정변수; 화폐 단위\\
$w^s$ & 시나리오 $s$의 페널티가 $\tau$를 초과하는 양 & 결정변수; 화폐 단위\\
$\mathcal R_\zeta(q)$ & 2단계 비용의 CVaR, 식 \eqref{eq:cvar} & $q$의 함수; 화폐 단위\\
$H(q,D^s)$ & 시나리오 $s$의 2단계(recourse) 비용, 식 \eqref{eq:second} & $q$의 함수; 화폐 단위\\
$B$ & 진입 예산 또는 진입 가능 도시 수 (변형 모형) & 입력; 화폐 단위 또는 도시 수\\
\end{longtable}}

\subsection{기본 표기}
{\footnotesize
\begin{longtable}{@{}L{0.10\textwidth}L{0.43\textwidth}L{0.41\textwidth}@{}}
\toprule 기호 & 의미 & 종류 및 단위\\\midrule\endfirsthead
\toprule 기호 & 의미 & 종류 및 단위\\\midrule\endhead\bottomrule\endlastfoot
$C$ & 후보 도시 집합 & 입력 집합; 도시\\
$S$ & 수요 시나리오 집합 & 입력 집합; 잔차 인덱스 집합, $|S|=N$\\
$y_c$ & 도시 $c$에 진입하면 1, 아니면 0 & 1단계 결정변수; 이진\\
$q_c$ & 도시 $c$에 배치하는 차량 대수 & 1단계 결정변수; 대 (정수)\\
$\kappa_c$ & 차량 한 대가 평가 기간 동안 처리할 수 있는 통행 수 & 척도 파라미터; 통행/대\\
$Q$ & 배치 가능한 총 차량 대수 (보유 차량; 취득비·감가상각은 매몰비용) & 입력; 대\\
$\bar q_c$ & 도시별 차량 배치 상한 & 입력; 대\\
$F_c$ & 도시 진입 및 기본 인프라 고정비 & 입력; 평가 기간 화폐 단위\\
$c_c^+$ & 잉여 용량 1단위당 손실. & 입력; 화폐/통행\\
$c_c^-$ & 미충족 잠재수요 1단위당 기회손실. 식 \eqref{eq:absorb} & 입력; 화폐/통행\\
$M_c$ & 도시 $c$의 1년간 \textbf{모든 수단을 합한 잠재 이동 기회} 수 & 입력; 이용 기회 수. 데이터 또는 가정\\
$p_c^s$ & 도시 $c$의 이동 기회가 로보택시를 선택할 확률 & 회귀로 생성; $[0,1]$. $p_c^s=\sigma(\hat\beta_0+\hat\beta_1x_c+\hat\varepsilon^s)$\\
$D_c^s$ & 로보택시 선택 수요 & 유도 입력; 통행. $D_c^s=M_c p_c^s$\\
$O_c^s$ & 잉여 서비스 용량 & 2단계 결과변수; 통행. $O_c^s=(\kappa_c q_c-D_c^s)^+$\\
$U_c^s$ & 미충족 로보택시 수요 & 2단계 결과변수; 통행. $U_c^s=(D_c^s-\kappa_c q_c)^+$\\
\end{longtable}}


%%%%%%%%%%%%%%%%%%%%%%%%%%%%%%%%%%%%%%%%
\section{수치 실험: 합성 데이터}

\subsection{전체 흐름}

참 수요 모형을 알고 있는 합성 환경에서 (i) 추정--최적화 파이프라인이 작동하는지, (ii) 참 모형이 있어야만 잴 수 있는 양(out-of-sample 비용, context의 가치, DRO의 효과)을 측정한다. 흐름은 다섯 단계이다.
\begin{enumerate}[nosep]
  \item 후보 도시 입력 $(x_c, M_c, \kappa_c, \bar q_c, F_c)$과 참 수요 모형 $p=\sigma(\beta_0+\beta_1 x+\varepsilon)$을 설정한다.
  \item 참 모형에서 훈련 자료 $(x^k,p^k)_{k=1}^N$을 생성하고 로짓 OLS로 $\hat f$와 잔차 $\hat\varepsilon^k$를 추정한다.
  \item 후보 도시별 시나리오 $D_c^s=M_c\,\sigma(\hat\beta_0+\hat\beta_1x_c+\hat\varepsilon^s)$를 만들어 확장형 \eqref{eq:ext}를 푼다.
  \item 참 모형에서 따로 생성한 평가 시나리오 1{,}000개로 해 $(y,q)$의 out-of-sample 기대비용과 꼬리비용을 계산한다.
  \item 설정 $(N, c^-, \zeta)$을 바꾸어 가며 2--4를 반복하고, context 없는 벤치마크 및 변형 \eqref{eq:budget}과 비교한다.
\end{enumerate}

\subsection{합성 데이터 생성}\label{sec:data}

\paragraph{후보 도시 입력.}
후보 도시는 Waymo가 공표한 향후 진출 후보 18개 도시이다. 각 도시를 중심 county 하나로 대응시키고(예: Chicago $\to$ Cook County), context $x_c$는 ACS 2019--2023 5년 추정 표 B08201의 무차량 가구 비율 실제값을 쓴다 (범위 $0.054$--$0.777$; New York County 0.777은 맨해튼만 포함한 극단값). 잠재 이동 기회 $M_c$는 같은 ACS 표의 총 가구 수에 비례상수 10을 곱해 두었으며, 이 비례상수는 toy용 가정이다 (실데이터에서는 NHTS $\times$ ACS 교차분류로 대체). 나머지 입력은 모든 도시에 공통으로 둔다 (표~\ref{tab:param}).

\paragraph{참 수요 모형.}
도시 $c$의 수요는 $D_c=M_c\,p_c$이고 불확실성은 선택확률 $p_c$에만 둔다. 참 모형은
\begin{equation}\label{eq:truth}
\logit(p)= \beta_0+\beta_1 x+\varepsilon,\qquad \varepsilon\sim N(0,\sigma_\varepsilon^2)
\end{equation}
이다. 로짓 형태를 쓰는 이유는 두 가지이다. 첫째, $p$는 ``이동 기회 하나가 로보택시를 선택할 확률''이므로 선택확률을 설명변수의 함수로 두는 표준형이 이산선택의 로짓이며, $p\in(0,1)$과 $x$에 대한 단조성이 자동으로 보장된다. 둘째, 관측된 로보택시 수요를 분석한 유일한 연구인 Loa et al.\ (2026)이 $\ln(\text{통행}+1)$을 종속변수로 한 로그-선형 공간시차 모형에서 무차량 가구 비율을 양(+)이고 유의한 변수(계수 $0.008$/\%p, $t=7.4$)로 보고하였고, $D=Mp$에서 로그-선형 통행 모형은 $p$가 작을 때 $\logit(p)\approx\ln p$와 같은 형태이다. 따라서 $\beta_1=0.8$은 그 계수를 0--1 척도로 환산한 값이다.

\paragraph{기본 파라미터.}
비용은 $c^+=1$로 정규화하고 $c^-$를 $1, 5, 10$으로 바꾼다(과잉 대 부족 비용 비율 $c^+:c^-$가 결정에 미치는 영향 확인 목적). 진입비는 $F_c=F_0\,c^-$로 두어 ``진입비 대 미충족 마진'' 비율이 $c^-$와 무관하게 유지되도록 하였다.

\begin{table}[H]\centering\small
\caption{기본 파라미터와 설정 근거}\label{tab:param}
\begin{tabular}{@{}llp{0.55\textwidth}@{}}
\toprule
파라미터 & 값 & 근거 \\ \midrule
$x_c$ & ACS 실제값 (0.054--0.777) & ACS 2019--2023 B08201, county 단위 \\
$M_c$ & 총 가구 수 $\times$ 10 & 가구 수는 ACS 실제값; 비례상수는 toy용 가정 \\
$\kappa_c$ & 8{,}760 통행/대/년 & 24통행/일 $\times$ 365일. Waymo 공개 수치(주당 유료 탑승 약 50만 건, 2026년 3월; 차량 3{,}000대 이상, 2025년 12월)에서 역산한 대당 일 24건 \\
$\bar q_c$ & 150 대 & 가정 (민감도 50--300) \\
$Q$ & $0.6\times\mathbb E[\sum_c D_c]/\kappa$ & 가정. 평균 수요의 60\%만 처리 가능하도록 하여 배분이 비자명 \\
$F_0$ & $10^5$ ($c^-=1$ 기준) & 가정 (민감도 $0.25\times10^5$-- $4\times10^5$) \\
$c^+:c^-$ & $1:1,\ 1:5,\ 1:10$ & 가정 \\
$\beta_0,\beta_1,\sigma_\varepsilon$ & $-3.5,\ 0.8,\ 0.3$ & $\beta_1$은 Loa et al.\ 환산값; $\beta_0$, $\sigma_\varepsilon$은 가정 \\
$N$ & 15, 200 & 도시--기간 관측 규모를 부족/충분 경우로 가정. 훈련 관측 수 $=$ 시나리오 수 $|S|$ \\
$\zeta$ & 0, 0.5, 0.9 & 0 = ER-SAA; 0.9 = 기본값 \\
평가 시나리오 & 1{,}000 & 임의 가정 \\
\bottomrule
\end{tabular}
\end{table}

\subsection{훈련 자료와 평가 자료}

\paragraph{훈련 자료.}
``로보택시가 이미 운영된 도시--기간에서 관측된 $(x,p)$''를 흉내 낸다. 후보 도시 $x_c$의 범위에서 임의의 $x^k$를 균등 추출하고, 참 모형~\eqref{eq:truth}에서 $\varepsilon^k$를 뽑아 $p^k$를 계산한다 ($k=1,\dots,N$). 모형은 참 계수를 모른 채 이 $N$개로 $\hat\beta_0,\hat\beta_1$과 잔차 $\hat\varepsilon^k$를 추정하고, 최적화는 그 추정치로 만든 시나리오만 본다. 실데이터에서는 이 단계가 $x$ = ACS 비율, $p$ = 로보택시 통행 수$/M$인 관측 표로 대체된다.

\paragraph{평가 자료.}
``결정을 내린 뒤 실제로 닥친 수요''를 흉내 낸다. 후보 도시의 실제 $x_c$에 참 계수를 넣고 훈련과 독립인 $\varepsilon$ 1{,}000개를 뽑아 $D^{\mathrm{eval}}_{c}$ (후보 도시 18개 $\times$ 1{,}000)를 만든다. 최적화는 이 자료를 보지 못하며, 해 $(y,q)$가 정해진 뒤 식~\eqref{eq:second}의 닫힌 형태로 시나리오별 비용을 계산해 평균(기대비용)과 상위 10\% 평균(꼬리비용, $\CVaR_{0.9}$)으로 채점한다. 평가 자료는 $N$과 무관하게 고정되어 있어 $N=15$와 $N=200$으로 내린 결정을 같은 잣대로 비교할 수 있다. 훈련 자료로 채점한 값(in-sample)은 최적화가 그 자료에 맞춰졌으므로 실제보다 낙관적이며, 표~\ref{tab:main}에서 in-sample과 OOS가 다른 이유이다.

\paragraph{풀이 방법 비교: 확장형 vs.\ L-shaped 분해.}
2단계 모형 \eqref{eq:first}--\eqref{eq:second}를 (i) 확장형 \eqref{eq:ext}로 통째로 푸는 방법과 (ii) master $(y,q,\tau)$와 시나리오별 subproblem으로 나누는 multi-cut L-shaped 분해로 각각 풀었다 ($N=15$, $c^-=5$, $\zeta=0.9$; 분해는 $\mathrm{UB}-\mathrm{LB}\le 10^{-2}\,|\mathrm{UB}|$, 즉 상대 허용오차 1\%에서 종료).

\begin{table}[H]\centering\small
\caption{확장형과 L-shaped 분해의 비교}\label{tab:solver}
\begin{tabular}{@{}llrrl@{}}
\toprule
설정 & 방법 & 목적값 & 시간 & 비고 \\ \midrule
$B=3$ & 확장형 & 16{,}036{,}515 & 0.4 s & 진입: Chicago, New York, San Diego \\
      & L-shaped & 16{,}140{,}861 & 29.5 s & 38회 반복, 진입 동일, 차이 0.65\% \\
$B$ 없음 & 확장형 & 17{,}304{,}198 & 1.5 s & 진입: Chicago, New York, Philadelphia, Seattle \\
      & L-shaped & 17{,}378{,}353 & 54.5 s & 38회 반복, 차이 0.43\%, 진입 도시 2곳 상이 \\
\bottomrule
\end{tabular}
\end{table}

L-shaped의 목적값은 확장형보다 $B=3$에서 0.65\%, $B$ 없음에서 0.43\% 높았고, 진입 도시는 $B=3$에서는 동일하였으나 $B$ 없음에서는 2곳이 달랐다. L-shaped가 돌려주는 값은 종료 시점의 실행가능해 비용(상계)이므로 이 차이는 허용오차에서 일찍 멈춘 결과이며, $B$ 없음에서 다른 도시 조합이 0.43\% 높은 비용으로 실행가능함을 보여준다. 계산시간은 확장형 0.4 s, 1.5 s에 대해 L-shaped 29.5 s, 54.5 s로 확장형이 훨씬 빨랐는데, 2단계 비용이 affine 두 개의 max라 확장형의 $O,U$ 변수가 cut을 처음부터 모두 넣은 것과 같기 때문이다. $F_c y_c$ 항을 제약 $B$로 옮기면($\sum_{c\in C}y_c\le B$) 분해의 수렴 시간이 줄어든다. 이후 실험은 모두 확장형으로 풀고, decision-dependent 확장 시 다시 비교한다.

\subsection{결과 분석}

실험 설정은 \S\ref{sec:data}와 같고, 참 모형에서 계산된 총 차량은 $Q=267$대이다. 모든 해는 참 모형의 시나리오 1{,}000개로 OOS 평가한다(기대비용과 꼬리비용 = 최악 10\% 평균). 아래 다섯 질문에 대해 답, 근거, 해당 표·그림 순으로 정리한다.

%---------------------------------------------------------------
\paragraph{Q1. 로짓 $p$ 회귀로 수요를 추정해도 되는가?}
표~\ref{tab:fit}에서 $N=200$이면 $\hat\beta_1=0.67$, $\hat\sigma_\varepsilon=0.30$으로 참값($0.8$, $0.3$) 근처를 회복한다. 잔차를 시나리오로 합쳐 쓰기 위한 조건(잔차가 $x$와 무관하게 평균 0, 분산 일정)은 그림~\ref{fig:diag}의 잔차 산점도로 확인된다. 검정 $R^2$가 $0.12$--$0.14$로 낮은 것은 설정상 $x$가 수요 변동의 약 13\%만 설명하기 때문이며, $x$ 없이 평균만 쓰면 $R^2\approx0$이므로 $x$의 설명력은 작지만 분명히 있다.

\begin{table}[H]\centering\small
\caption{로짓 OLS 추정 결과 (오차 지표는 $p$ 척도, 검정 자료 1{,}000개)}\label{tab:fit}
\begin{tabular}{@{}rrrrrrrrr@{}}
\toprule
$N$ & $\hat\beta_0$ & $\hat\beta_1$ & $\hat\sigma_\varepsilon$ & test MSE & test MAE & test ME & test $R^2$ & no-context $R^2$ \\ \midrule
15  & $-3.655$ & $1.114$ & $0.325$ & $0.0002$ & $0.0101$ & $0.0461$ & $0.124$ & $-0.011$ \\
200 & $-3.479$ & $0.665$ & $0.298$ & $0.0002$ & $0.0099$ & $0.0479$ & $0.143$ & $-0.013$ \\
참값 & $-3.500$ & $0.800$ & $0.300$ & & & & & \\
\bottomrule
\end{tabular}
\end{table}

\begin{figure}[H]\centering
\includegraphics[width=\textwidth]{fig1_f_diagnostics.png}
\caption{$f$ 추정 진단. (a) 로짓 척도의 자료·적합선·참 직선, (b) 잔차 분포, (c) 잔차 대 $x$, (d) 검정 자료의 예측 $p$ 대 실제 $p$. 윗줄 $N=15$, 아랫줄 $N=200$.}\label{fig:diag}
\end{figure}

%---------------------------------------------------------------
\paragraph{Q2. $c^+:c^-$ 비율은 무엇을 쓰는가?}
비율은 자료가 없어 선택하는 입력이며, 실험의 역할은 결론이 비율에 따라 바뀌지 않음을 보이는 것이다. 표~\ref{tab:main}에서 $c^-=1,5,10$ 모두 (i) context 모형이 context 없는 모형보다 OOS 기대비용이 낮고, (ii) $\zeta$가 커질수록 진입 도시 수가 줄며, (iii) Chicago와 New York은 18개 설정 모두에서, San Diego는 17개 설정에서 선택된다. 그림~\ref{fig:sens}의 $c^-$ 패널에서 비용은 $c^-$에 비례해 커지지만($c^-=1,2,5,10,20$에서 3.08, 5.79, 14.22, 27.63, 54.45M) 세 모형의 상대 순서는 바뀌지 않는다. 따라서 기준 설정은 미충족이 잉여보다 비싼 $1:5$로 두고 나머지 비율은 보고만 한다.

\begin{table}[H]\centering\footnotesize
\caption{설정별 최적 진입 도시 수, in-sample 목적값, OOS 비용 ($c^+=1$ 단위) 및 context 이득}\label{tab:main}
\begin{tabular}{@{}rrr r rrr rr r@{}}
\toprule
$N$ & $c^-$ & $\zeta$ & 진입 수 & in-sample & OOS 평균 & OOS 꼬리 & no-ctx 평균 & no-ctx 꼬리 & 이득 \\ \midrule
15  & 1  & 0   & 9 & 2{,}490{,}926 & 2{,}695{,}476 & 4{,}857{,}473 & 2{,}758{,}361 & 4{,}657{,}473 & 2\% \\
    &    & 0.5 & 6 & 2{,}864{,}469 & 2{,}724{,}601 & 4{,}557{,}473 & 3{,}112{,}997 & 4{,}357{,}473 & 12\% \\
    &    & 0.9 & 4 & 3{,}460{,}840 & 3{,}147{,}884 & 4{,}410{,}069 & 3{,}869{,}252 & 4{,}739{,}552 & 19\% \\
    & 5  & 0   & 6 & 11{,}386{,}313 & 12{,}541{,}832 & 22{,}787{,}364 & 12{,}702{,}986 & 22{,}787{,}364 & 1\% \\
    &    & 0.5 & 5 & 13{,}635{,}159 & 12{,}712{,}913 & 22{,}287{,}364 & 14{,}896{,}198 & 22{,}208{,}291 & 15\% \\
    &    & 0.9 & 4 & 17{,}304{,}198 & 13{,}787{,}444 & 21{,}975{,}797 & 15{,}636{,}705 & 23{,}199{,}859 & 12\% \\
    & 10 & 0   & 6 & 22{,}412{,}124 & 24{,}728{,}335 & 45{,}574{,}729 & 25{,}166{,}440 & 45{,}574{,}729 & 2\% \\
    &    & 0.5 & 5 & 27{,}186{,}803 & 25{,}027{,}407 & 44{,}574{,}729 & 28{,}863{,}429 & 44{,}263{,}095 & 13\% \\
    &    & 0.9 & 4 & 34{,}608{,}397 & 27{,}416{,}340 & 43{,}920{,}442 & 30{,}089{,}296 & 45{,}739{,}579 & 9\% \\ \midrule
200 & 1  & 0   & 8 & 2{,}599{,}840 & 2{,}662{,}764 & 4{,}757{,}473 & 2{,}748{,}503 & 4{,}657{,}473 & 3\% \\
    &    & 0.5 & 6 & 3{,}054{,}078 & 2{,}712{,}652 & 4{,}557{,}473 & 2{,}980{,}460 & 4{,}457{,}473 & 9\% \\
    &    & 0.9 & 4 & 4{,}415{,}710 & 3{,}084{,}216 & 4{,}357{,}473 & 3{,}727{,}040 & 4{,}610{,}719 & 17\% \\
    & 5  & 0   & 6 & 12{,}225{,}772 & 12{,}463{,}377 & 22{,}787{,}364 & 12{,}722{,}661 & 22{,}787{,}364 & 2\% \\
    &    & 0.5 & 5 & 14{,}870{,}099 & 12{,}674{,}534 & 22{,}287{,}364 & 13{,}031{,}319 & 22{,}287{,}364 & 3\% \\
    &    & 0.9 & 3 & 22{,}045{,}780 & 14{,}222{,}971 & 21{,}617{,}923 & 14{,}901{,}907 & 22{,}265{,}378 & 5\% \\
    & 10 & 0   & 6 & 24{,}168{,}288 & 24{,}654{,}132 & 45{,}574{,}729 & 25{,}128{,}612 & 45{,}574{,}729 & 2\% \\
    &    & 0.5 & 5 & 29{,}646{,}948 & 24{,}957{,}044 & 44{,}574{,}729 & 25{,}611{,}150 & 44{,}574{,}729 & 3\% \\
    &    & 0.9 & 3 & 44{,}013{,}689 & 27{,}629{,}179 & 43{,}180{,}752 & 28{,}889{,}111 & 44{,}519{,}917 & 4\% \\
\bottomrule
\end{tabular}
\end{table}

%---------------------------------------------------------------
\paragraph{Q3. 왜 contextual ER-DRO인가, 조건은 만족하는가?}
조건은 Q1에서 확인되었다(잔차의 $x$ 독립성과 $\hat f$의 일치성). 효과는 두 부분이다. \emph{Context의 효과}: 표~\ref{tab:main}에서 context 모형의 OOS 기대비용 이득은 $\zeta=0$에서 1--3\%, $\zeta=0.5$에서 3--15\%, $\zeta=0.9$에서 4--19\%이다. 이득이 작은 것은 $\beta_1=0.8$에서 $x$의 설명력이 작기 때문이며, 그림~\ref{fig:sens}의 $\beta_1$ 패널에서 $\beta_1=0$이면 두 모형의 비용이 같고 $\beta_1=0.8, 1.5, 3, 5$에서 이득은 1\%, 16\%, 25\%, 41\%로 커진다. \emph{DRO의 효과}: $\zeta$는 기대비용과 꼬리비용을 맞바꾼다. $N=200$, $c^-=5$에서 $\zeta=0\to0.5\to0.9$일 때 기대비용은 $12.46\to12.67\to14.22$M($+1.7\%$, $+14\%$), 꼬리비용은 $22.79\to22.29\to21.62$M($-2.2\%$, $-5.1\%$)이다(그림~\ref{fig:zeta}). 한편 context 없는 모형은 $N=15$, $\zeta=0.9$에서 꼬리비용이 오히려 $\zeta=0$보다 커진다($c^-=5$: 22.79M$\to$23.20M; $c^-=1,10$도 동일). $x$ 없이 잔차 15개만으로 꼬리를 맞추려다 생기는 현상이며, DRO는 context와 함께 써야 함을 뜻한다.

\begin{figure}[H]\centering
\includegraphics[width=0.9\textwidth]{toy_oos_vs_zeta.png}
\caption{$\zeta$에 따른 OOS 기대비용(좌)과 꼬리비용(우), $c^+:c^-=1:5$.}\label{fig:zeta}
\end{figure}

%---------------------------------------------------------------
\paragraph{Q4. 자료가 적어도 이 모형을 써도 되는가?}
context를 쓰면 된다. 표~\ref{tab:main}의 $\zeta=0.9$ 행에서 $N=15$와 $N=200$의 OOS 기대비용 차이는 $c^-=1,5,10$에서 2.1\%, 3.1\%, 0.8\%, 꼬리비용 차이는 1.2--1.7\%이다. $\hat\beta_1$이 $N=15$에서 39\% 과대추정되고 $\zeta=0.9$에서는 상위 2개 시나리오만 목적함수에 들어가는데도 그렇다. 그림~\ref{fig:sens}의 $N$ 패널에서 ER-SAA(context)는 $N=10$부터 $200$까지 1.4\% 안에서 평평하고, ER-DRO(context)는 13.19--14.24M, context 없는 모형은 14.38--16.14M으로 흔들린다. 즉 자료가 적을수록 context 모형의 이점이 커진다.

%---------------------------------------------------------------
\paragraph{Q5. 민감도.}
그림~\ref{fig:sens}는 기준 설정($N=200$, $c^-=5$, $\zeta=0.9$)에서 파라미터를 하나씩 바꾼 결과이다. 진입비 $F_0$와 변동성 $\sigma_\varepsilon$이 오르면 비용이 오르고($F_0$: 12.04$\to$18.73M), 총 차량 $Q$와 대당 처리량 $\kappa$가 오르면 비용이 내린다($Q$ 비율 0.2$\to$1.0: 16.79$\to$11.45M). $Q$가 평균 수요의 20\%로 작으면 세 모형의 비용이 거의 같아지는데, 차량이 어떻게 배분해도 전부 쓰이기 때문이다. 도시별 상한 $\bar q_c$는 ER-SAA에서는 100 이상, ER-DRO에서는 150 이상에서 결과를 바꾸지 않는다. 표~\ref{tab:budget}의 변형 \eqref{eq:budget}에서는 진입 도시 수 $B=2,3,5$로 늘릴수록 2단계 비용이 내려가고($N=200$, $c^-=5$: 14.57$\to$12.73$\to$11.74M), 처음 세 도시는 $N$과 $c^-$에 무관하게 Chicago, New York $\to$ San Diego 순으로 안정적이다. $B=5$의 구성은 $N$에 따라 달라지며, $N=200$, $c^-=5,10$에서는 New York 대신 Philadelphia, Pittsburgh, Tampa가 들어온다.

\begin{figure}[H]\centering
\includegraphics[width=\textwidth]{fig2_sensitivity.png}
\caption{민감도 분석. 기준 $N=200$, $c^-=5$, $c^+=1$, $\zeta=0.9$에서 파라미터를 하나씩 변경. 선: context ER-DRO ($\zeta=0.9$), context ER-SAA ($\zeta=0$), context 없음 ($\zeta=0.9$).}\label{fig:sens}
\end{figure}

\begin{table}[H]\centering\small
\caption{변형 \eqref{eq:budget}: 진입 도시 수 $B$에 따른 OOS 2단계 비용과 선택 도시 ($\zeta=0.9$)}\label{tab:budget}
\begin{tabular}{@{}rrrrl@{}}
\toprule
$N$ & $c^-$ & $B$ & OOS 2단계 비용 & 진입 도시 \\ \midrule
15  & 1  & 2 & 3{,}315{,}579 & Chicago, New York \\
    &    & 3 & 3{,}084{,}918 & + San Diego \\
    &    & 5 & 2{,}358{,}358 & + Philadelphia, Seattle \\
    & 5  & 2 & 14{,}339{,}887 & Chicago, New York \\
    &    & 3 & 12{,}668{,}561 & + San Diego \\
    &    & 5 & 10{,}305{,}798 & + Philadelphia, Seattle \\
    & 10 & 2 & 28{,}111{,}826 & Chicago, New York \\
    &    & 3 & 24{,}631{,}627 & + San Diego \\
    &    & 5 & 20{,}068{,}402 & + Philadelphia, Seattle \\ \midrule
200 & 1  & 2 & 3{,}395{,}462 & Chicago, New York \\
    &    & 3 & 3{,}198{,}004 & + San Diego \\
    &    & 5 & 2{,}440{,}092 & + Detroit, Philadelphia \\
    & 5  & 2 & 14{,}572{,}482 & Chicago, New York \\
    &    & 3 & 12{,}726{,}413 & + San Diego \\
    &    & 5 & 11{,}735{,}447 & Chicago, Philadelphia, Pittsburgh, San Diego, Tampa \\
    & 10 & 2 & 28{,}274{,}685 & Chicago, New York \\
    &    & 3 & 24{,}625{,}772 & + San Diego \\
    &    & 5 & 22{,}821{,}183 & Chicago, Philadelphia, Pittsburgh, San Diego, Tampa \\
\bottomrule
\end{tabular}
\end{table}

%---------------------------------------------------------------
\paragraph{요약.}
(Q1) 로짓 $p$ 회귀는 참 계수를 회복하고 잔차 합동 조건을 만족한다. (Q2) 결론은 $c^+:c^-$ 비율에 둔감하므로 $1:5$를 기준으로 둔다. (Q3) context는 $x$가 수요를 좌우하는 정도만큼 비용을 줄이고(현 설정 1--3\%, 강한 context 25--41\%), DRO는 기대비용을 내주고 꼬리비용을 줄이며, 둘은 함께 써야 한다. (Q4) 관측 15개로도 비용이 3\% 안에서 유지된다. (Q5) 비용은 $Q$, $\kappa$, $F_0$, $\sigma_\varepsilon$의 규모에 민감하고, Chicago와 New York은 모든 설정에서, San Diego는 17/18 설정에서 선택된다.
\\
\\
==========================================================

\section{실데이터 적용 계획}

\subsection{기호별 사용 자료}

\begin{table}[H]\centering\small
\caption{모형 입력과 실데이터 출처}\label{tab:realdata}
\begin{tabular}{@{}L{0.17\textwidth}L{0.30\textwidth}L{0.33\textwidth}L{0.12\textwidth}@{}}
\toprule
입력 & 실데이터 & 출처 & 단위 \\ \midrule
$p^k$ (훈련 종속변수) & 로보택시 통행 수 $\div M^k$ & CPUC AV Passenger Service 분기 보고서 (사업자 제출; 캘리포니아) & tract $\times$ 기간 \\
$x^k$ (훈련 context) & 무차량 가구 비율 & ACS 5년 추정 B08201 (\texttt{\_002E/\_001E}) & tract \\
$M^k$ (훈련) & $\sum_g$ (차량 $g$대 보유 가구 수)$\times$(NHTS $g$집단 통행률)$\times$일수 & ACS B08201 (\texttt{\_002E}--\texttt{\_006E}) + NHTS 2022 & tract $\times$ 기간 \\
$x_{cj}$ (후보 도시 tract) & 무차량 가구 비율 & ACS B08201, 후보 county 안의 tract & tract \\
$M_{cj}$ (후보 도시 tract) & 위와 같은 공식, 연 단위 & ACS + NHTS & tract $\times$ 년 \\
$\kappa$ & 8{,}760 통행/대/년 & Waymo 공개 수치 역산 & -- \\
$Q$, $\bar q_c$, $F_0$ & toy와 같은 가정 & -- & -- \\
$c^+:c^-$ & $1{:}1,\ 1{:}5,\ 1{:}10$ & 비용 자료 없음 $\to$ 비율 & -- \\
$\zeta$ & 0.9 고정 & -- & -- \\
\bottomrule
\end{tabular}
\end{table}

\subsection{절차}

운영 도시(샌프란시스코 등)의 tract $k=1,\dots,N$에서 $x^k$, $M^k$, 통행 수 $T^k$를 관측하고 $p^k=T^k/M^k$로 둔다. 후보 도시 $c$의 tract $j\in J_c$에는 $x_{cj}$, $M_{cj}$만 있다(통행 없음).

\paragraph{1. 회귀 (운영 도시 tract).}
\begin{equation}
\logit(p^k)=\beta_0+\beta_1x^k+\varepsilon^k,\qquad
(\hat\beta_0,\hat\beta_1)=\arg\min_{\beta}\sum_{k=1}^N\big(\logit(p^k)-\beta_0-\beta_1x^k\big)^2,\qquad
\hat\varepsilon^k=\logit(p^k)-\hat\beta_0-\hat\beta_1x^k.
\end{equation}

\paragraph{2. 시나리오 (후보 도시 tract에 적용; 잔차 하나 $=$ 시나리오 하나).}
\begin{equation}
p_{cj}^{\,s}=\sigma\big(\hat\beta_0+\hat\beta_1x_{cj}+\hat\varepsilon^{s}\big),\qquad
D_c^{\,s}=\sum_{j\in J_c}M_{cj}\,p_{cj}^{\,s},\qquad s=1,\dots,N.
\end{equation}
$c$는 후보도시, $j$는 그 도시 안의 tract(수요 계산 단위).

\paragraph{3. 최적화.} $D_c^s$를 확장형 \eqref{eq:ext}에 넣어 $(y_c,q_c)$를 구한다. 결정은 도시 단위, 수요는 tract 합이다.

\paragraph{4. 도시 간 교차검증.} CPUC 자료에 San Francisco외 LA 도 있다. 샌프란시스코 tract들로 $\hat f$ 추정하고 LA tract의 x를 넣어 예측한 p를 LA의 실제 p와 비교한다.(MAE, $R^2$, 잔차 평균 0인지, 퍼짐이 SF 잔차와 같은지) 만약 LA 잔차 평균이 0에서 크게 벗어나면 외삽에 한계가 있다는 것이다.


\subsection{전제 조건}

\begin{enumerate}[nosep]
  \item \textbf{같은 관계(외삽)}: 운영 도시 tract에서 추정한 $\beta_0,\beta_1$이 후보 도시 tract에도 그대로 성립한다.
  \item \textbf{잔차 합동(pooling)}: $\varepsilon$의 분포가 tract와 $x$에 무관하다(평균 0, 등분산). 그래야 운영 도시의 잔차를 후보 도시 어느 tract에나 붙일 수 있다.
  \item \textbf{같은 단위}: 추정과 적용이 같은 공간 단위(tract)이고, $p$의 분자·분모가 같은 기간이며, $M$이 두 도시에서 같은 공식으로 계산된다.
  \item \textbf{범위 안}: 후보 tract의 $x_{cj}$가 훈련 $x^k$의 범위 안에 있다.
  \item \textbf{수요 외생성}: $p$가 결정 $q_c$에 의존하지 않는다(decision-independent; 이후 확장에서 완화).
\end{enumerate}

\subsection{작은 숫자 예시}

샌프란시스코(운영 중)에 tract 3개, 시카고(후보)에 tract 2개만 있다고 하자. 기간은 1년이다.

\paragraph{1단계. 훈련 자료 (샌프란시스코 tract).}
\begin{table}[H]\centering\small
\begin{tabular}{@{}lrrrrrr@{}}
\toprule
SF tract & 가구 수 & 무차량 가구 & $x$ & $M$ (가구$\times$10) & 로보택시 통행 (CPUC) & $p=$ 통행$/M$ \\ \midrule
A & 1{,}000 & 100 & 0.10 & 10{,}000 & 300 & 0.03 \\
B & 1{,}000 & 300 & 0.30 & 10{,}000 & 600 & 0.06 \\
C & 1{,}000 & 500 & 0.50 & 10{,}000 & 800 & 0.08 \\
\bottomrule
\end{tabular}
\end{table}
$x$와 가구 수는 ACS, 통행 수는 CPUC에서 온다. 관측은 $N=3$개이다. 도시 단위로 하면 샌프란시스코 전체가 한 줄($N=1$)이 되어 회귀가 불가능하므로 tract로 내려간다.

\paragraph{2단계. 회귀.} $\logit(p)$를 $x$에 적합하면 $\hat\beta_0=-3.67$, $\hat\beta_1=2.58$, 잔차 $\hat\varepsilon=(-0.07,\,+0.14,\,-0.07)$이다. 이 세 값이 그대로 시나리오 3개가 된다.

\paragraph{3단계. 후보 도시에 적용.} 시카고에는 통행 자료가 없고 ACS만 있다.
\begin{table}[H]\centering\small
\begin{tabular}{@{}lrr@{}}
\toprule
Chicago tract & $x$ & $M$ \\ \midrule
D & 0.20 & 20{,}000 \\
E & 0.60 & 5{,}000 \\
\bottomrule
\end{tabular}
\end{table}
시나리오 1(잔차 $-0.07$): tract D의 $p=\sigma(-3.67+2.58\times0.20-0.07)=0.0385$, tract E의 $p=\sigma(-3.67+2.58\times0.60-0.07)=0.1012$, 수요 $D=20{,}000(M_D)\times0.0385(p_D)+5{,}000(M_E)\times0.1012(p_E)=770+506=1{,}276$. 같은 식으로 시나리오 2(잔차 $+0.14$)는 $940+609=1{,}548$, 시나리오 3은 $1{,}276$이다. 이 세 값이 최적화의 $D_{\text{Chicago}}^{s}$, $s=1,2,3$이다. 결정 $(y,q)$는 도시 하나에 대해 내리고, 수요만 tract 합으로 계산한다.


\paragraph{시간 단위.} 위 표가 1년이 아니라 1개월이어도 $p=$ 통행$/M$은 같은 척도이므로 $\hat f$를 그대로 쓸 수 있고, 시카고 $M$을 연간으로 두면 연간 수요가 나온다. 다만 월 자료는 잔차가 더 퍼져 있어 시나리오의 변동이 커진다.

%\paragraph{toy 코드와의 대응.} 1단계 표가 \texttt{make\_training\_data}의 출력, 2단계가 \texttt{fit\_logit\_ols}, 3단계가 \texttt{make\_scenarios}이다. 실데이터 단계에서는 도시당 $x_c$ 하나 대신 tract 벡터 $x_{cj}$를 넣고 tract 합을 취하도록 \texttt{make\_cities}와 \texttt{make\_scenarios}만 바꾸면 된다.

\end{document}